%% =====================================================================
%%  3 Assumptions and Their Consequences
%% =====================================================================
\section{Assumptions and Their Consequences}
\label{sec:assumptions}

Every answer in this paper is conditional on the eight assumptions below. Each
is stated, justified against the literature or the problem data, and followed
by the change in the answer if it fails. Where the change can be measured by
rerunning an estimator with the effect altered, the measured value is quoted;
where it requires a separate estimate, the estimate is given in one line.
A1--A5 concern the flight model, A6--A7 the data, and A8 the definition of a
goal. Several of these assumptions are sloppy in the technical sense: the
prediction depends on combinations of parameters that the data leave almost
free, so a statement about a parameter is not by itself a statement about the
prediction \citep{gutenkunst2007}.

\begin{enumerate}[label=\textbf{A\arabic*},leftmargin=*,itemsep=6pt]

\item \textbf{Rigid, non-deforming ball; no contact with post or net.}
The model needs a constant radius and cross-section, and it stops at the goal
plane, before the net. The Laws fix the circumference at \SIrange{68}{70}{cm}
and the mass at \SIrange{410}{450}{g} \citep{ifab2025}, and the problem's
$z_0=\SI{0.11}{m}$ gives $R=\SI{0.110}{m}$. At \SI{31}{m/s} the aerodynamic
pressure $\tfrac12\rho v^2\approx\SI{0.6}{kPa}$ is of order one per cent of
the \SIrange{0.6}{1.1}{atm} internal pressure, so the ball stays spherical to
well under a millimetre. A \SI{2}{mm} error in the effective radius changes
the cross-section by $2\Delta R/R\approx3.6$ per cent and the Magnus force,
which carries a further factor $R$, by $3\Delta R/R\approx5.5$ per cent;
applied to the \SI{0.67}{m} of entry the transverse forces produce, this is a
few centimetres. Contact with the post or the net lies outside the model: it
would move the crossing by at least the width of the post, and such
trajectories are not described by the free-flight input.

\item \textbf{Constant drag coefficient during the flight.}
The speed falls from \SI{32.6}{m/s} at the kick to roughly \SI{25}{m/s} at
the plane, a quarter of its value, so the Reynolds number stays above the drag
crisis at $\mathrm{Re}\approx9\times10^4$--$1.3\times10^5$, where
$C_D(\mathrm{Re})$ is flat at $0.17$--$0.22$ \citep{carre2005,goff2012}; that
flatness, not the particular value, is what licenses a single $C_D$
\citep{mehta1985}. Rather than fixing $C_D$, we hold
it in a Gaussian prior centred at $0.25$ with standard deviation $0.15$, which
covers $0.10$--$0.40$ at one standard deviation and contains both the
$0.17$--$0.22$ plateau and the $0.21$ measured on a spinning ball
\citep{asai2020}; the data then move all fits of this family into
$C_D=0.32$--$0.38$, above the plateau; that disagreement is examined below. Shifting the prior median across the plateau moves the
entry by about a centimetre, and the reason is structural: drag is
antiparallel to the velocity, so at fixed launch velocity it cannot change
$v_y/v_x$, and the lateral entry is $y=\SI{3.217}{m}$ for every $C_D$. Drag
does change the flight time, and through it the height: $z=\SI{1.957}{m}$ with
gravity alone against \SI{1.879}{m} with drag, a shift of \SI{8}{cm}.

\item \textbf{Transverse force linear in spin, fixed axis, exponential decay.}
The force is
$\vect{F}_M=\tfrac12\rho A C_{L0}R(\vect{\omega}\times\vect{v})$ with
$C_L=C_{L0}S_p$, $C_{L0}=1.0$, a constant spin axis, and
$\omega(t)=\omega_0\ee^{-t/\tau}$ with $\tau=\SI{6}{s}$. This is the smallest
closure that leaves the three spin components free and keeps the force smooth
at zero spin. The linear law is supported by measurement: $C_L=0.21$, $0.24$,
$0.29$ at $S_p=0.18$, $0.22$, $0.34$ gives $C_L/S_p=1.17$, $1.09$, $0.85$, so
the adopted coefficient lies inside the measured range \citep{asai2020}. The
main line does not need the spin: it fits the transverse force directly, and a
unit lift slope turns its fitted magnitude into an implied spin of
\SI{9.6}{rev/s}, inside the measured \SIrange{2}{10}{rev/s}. The earlier
spinning-ball control, which did carry a spin vector and an exponential decay,
found $\SI{52 \pm 22}{rad/s}$ ($8.3\pm3.5$~rev/s) with $\tau=\SI{6}{s}$; varying
$\tau$ over \SIrange{2}{12}{s} there moved the entry by less than \SI{3}{cm},
and varying its lift coefficient over $0.7$--$1.3$ moved it by \SI{15}{cm}. So
the magnitude of the transverse force is what the data constrain; the
mechanism that produces it is not. That is why the six-parameter fit models the
force through two coefficients with a zero-centred prior of standard deviation
\SI{0.008}{m^{-1}} rather than through a spin vector, and why a non-spinning
ball's wake asymmetry \citep{texier2016} is a legitimate member of the same
family.

\item \textbf{Still air.}
No wind term is included, because two fixes cannot separate a uniform wind
from the launch velocity: the wind enters the drag as a relative velocity and
is absorbed, to first order, by a change of $\vect{v}_0$. Refitting with a
\SI{3}{m/s} crosswind and the launch velocity free moves the entry by
\SI{0.007}{m}, so wind of this strength is inside the model envelope rather
than a leading term. At fixed launch velocity a \SI{1}{m/s} crosswind
would displace the entry by about \SI{3}{cm}, but the refit shows that the
launch velocity absorbs almost all of it. A neglected wind is also model discrepancy rather
than parameter uncertainty, and discrepancy can bias a prediction without
widening its stated interval \citep{white2017}, so it is reported as a
systematic and not folded into $\sigma_{\rm pos}$.

\item \textbf{Ball radius \SI{0.110}{m}, mass \SI{0.430}{kg}.}
The model converts aerodynamic coefficients into accelerations through $R$ and
$m$. The Laws bracket the mass at \SIrange{410}{450}{g} and the circumference
at \SIrange{68}{70}{cm}, i.e.\ $R=\SIrange{0.108}{0.111}{m}$
\citep{ifab2025}. The problem's kick height selects $R=\SI{0.110}{m}$, and we
use $m=\SI{0.430}{kg}$, the middle of the legal band. Both drag and Magnus
accelerations scale as $1/m$, so the $\pm20$~g band, $\pm4.7$ per cent, scales
the transverse part of the entry by the same fraction; at \SI{0.67}{m} that
is about \SI{3}{cm}, $0.15$ of the \SI{0.20}{m} predictive standard deviation.
The radius enters more weakly because it also enters the area: the
$\pm\SI{2}{mm}$ band changes the accelerations by the few per cent estimated
in A1 and the entry by a few centimetres.

\item \textbf{Gaussian, independent, isotropic position noise with
$\sigma_{\rm pos}=\SI{5}{mm}$.}
The likelihood needs a noise model, and two fixes cannot estimate one. The
coordinates are printed to \SI{0.01}{m}, so half of the last printed digit is
the only defensible scale; this is not a camera calibration, and the
$\pm\SI{1.5}{cm}$ of goal-line technology \citep{fifa2020glt} bounds a
certified installation rather than these data. Open tracking is coarser: local
cubic fits over nine-frame windows leave residuals of \SI{0.056}{m} and
\SI{0.054}{m} in $x$ and $y$ for IDSSE and \SI{0.065}{m} and
\SI{0.064}{m} for Metrica, so a goal-line fix is better than a centimetre and
a broadcast feed worse, and the printed precision lies between them. The
consequence scales as a second difference: the acceleration uncertainty
$\sqrt6\,\sigma_{\rm pos}/\Delta t^2$ is \SI{4.9}{m/s^2} at \SI{5}{mm} and
\SI{19.6}{m/s^2} at \SI{20}{mm}, where it exceeds the \SI{16}{m/s^2} signal
itself, and across the twenty-five single-factor variants, four joint-stress cases and
an independent replication, P(clean goal) ranges from $0.735$ at \SI{20}{mm} to
$0.9997$ at \SI{2.5}{mm}, while the scenario means stay within
$y=\SIrange{2.687}{3.130}{m}$ and inside the model envelope. The noise level
is the largest single factor after the force law, and the ranking of the
estimators survives the whole range.

\item \textbf{The fixes are the ball centre, with no common offset.}
The inversion differences positions, so a common offset is a systematic, not a
random, error, and must be stated separately. The problem's
$z_0=\SI{0.11}{m}$ shows that the coordinates refer to the centre. Because the
acceleration is a second difference, a common offset $b$ in both fixes appears
in the constant-acceleration fit as
$2b(t_2-t_1)/[t_1t_2(t_1-t_2)]\approx-400b$ (with $b$ in metres): a
\SI{1}{cm} offset would be read as \SI{4}{m/s^2}, a quarter of the
\SI{16}{m/s^2} that the data imply, and would move the entry by \SI{0.21}{m},
four-tenths of the \SI{0.54}{m} width of the literature-prior control. Below
\SI{2.5}{mm} the effect is under \SI{5}{cm}. It is nevertheless
the assumption most worth testing with a third fix or a known launch point,
and it motivates the systematic line in the error budget.

\item \textbf{Crossing of the ball centre; whole-ball scoring.}
The crossing time and point are defined by the ball centre reaching
$x=\SI{11.0}{m}$, and a goal is scored when the whole ball has passed inside
the frame, so the clean-goal region is
$|y|\le 3.66-0.11=\SI{3.55}{m}$ and $0.11\le z\le 2.44-0.11=\SI{2.33}{m}$
\citep{ifab2025}. All probabilities quoted here use this whole-ball
convention. The distinction is not cosmetic: it shrinks the effective target
from $7.32~\text{m}\times2.44$~m to $7.10~\text{m}\times2.22$~m, a loss of
\SI{12}{\percent} of its area. With the nominal predictive standard deviations
$\sigma_y\approx\sigma_z\approx\SI{0.20}{m}$, the predicted centre lies
$3.7\sigma$ inside the post and $1.9\sigma$ below the crossbar; removing one
radius from each side lowers $P(\text{clean goal})$ from $0.993$ under a
centre-only convention to $0.971$, a difference of two percentage points that
comes almost entirely from the crossbar. A centre-only convention is not
directly comparable with the numbers reported here.

\end{enumerate}
